\subsection{模型一的建立与求解：无噪声时间配准与10 Hz轨迹重建}\label{sec:model-q1}

\subsubsection{模型一的建立}\label{sec:q1-establishment}

\noindent\textbf{（1）模型目标、输入输出与证据边界。}

附件 1 给出同一机器人轨迹的两组无噪声采样：方式 1 为 4 Hz，方式 2 为 5
Hz。两台设备开机先后不同，故设备时间轴之间存在固定偏差。本问需要估计该
偏差，并在不外推的前提下生成 10 Hz 二维位置轨迹。

数据审计表明，方式 1、方式 2 分别有 3000、4001 个有效观测，采样间隔严格
为 0.25 s、0.20
s，无缺失和重复时间戳。由于题面明确无噪声，本问不进行滤波、
平滑或异常值删除。

\noindent\textbf{（2）方法选择与无噪声边界。}

方案 A 直接在粗网格上最小化位置误差，简单但最终精度受网格步长限制；方案
B
先用异频统一网格上的速度互相关产生候选，再以全局位置误差扫描和有界连续优化
精估时间偏差，可同时防止周期轨迹的伪相关峰；方案 C
采用动态时间规整允许非线性
时间伸缩，但题面只存在固定开机偏差，额外自由度会改变问题含义。因此选择方案
B。

连续轨迹比较 Linear、Cubic Spline、PCHIP 和 Akima
四种插值。最终以对齐后的 跨传感器位置 RMSE
为主判据，并兼顾无噪声轨迹的连续可导性，选择 Cubic Spline；其中 PCHIP 的保形性质及其使用边界参照分段三次插值理论\cite{fritsch1980}。

\noindent\textbf{（3）时间映射与位置误差目标。}

记两种定位数据为

\begin{equation}
D_m=\{(t_i^{(m)},x_i^{(m)},y_i^{(m)})\},\qquad m=1,2,
\label{eq:q1-data}
\end{equation}

并分别建立连续二维轨迹 \(\mathbf r_m(t)\)。定义方式 2 的校正时间为

\begin{equation}
t_{2,\mathrm{aligned}}=t_2+\Delta t.
\label{eq:q1-time-map}
\end{equation}

因此 \(\Delta t<0\) 表示方式 2 的时间戳需要整体向前平移。对候选偏差
\(\Delta t\)，公共有效区间为

\begin{equation}
\Omega(\Delta t)=
\left[
\max(t_{1,\min},t_{2,\min}+\Delta t),
\min(t_{1,\max},t_{2,\max}+\Delta t)
\right].
\label{eq:q1-overlap}
\end{equation}

在公共评价时刻 \(\{t_k\}\) 上定义位置均方误差

\begin{equation}
J(\Delta t)=\frac{1}{N}
\sum_{k=1}^{N}
\left\|\mathbf r_1(t_k)-\mathbf r_2(t_k-\Delta t)\right\|_2^2,
\label{eq:q1-objective}
\end{equation}

从而得到

\begin{equation}
\Delta t^*=\arg\min_{\Delta t}J(\Delta t).
\label{eq:q1-optimum}
\end{equation}

\subsubsection{模型一的求解}\label{sec:q1-solution}

求解时先在 0.1 s 公共网格上计算二维速度特征互相关，保留多个相关峰；随后在
全部可行偏差范围内以 0.5 s 步长扫描位置
MSE，最后在最优候选附近用有界标量
优化精细求解。精细优化使用固定评价网格，避免候选偏差改变公共区间起点所造成的
数值非光滑。

\noindent\textbf{（1）同时刻观测合并与10 Hz重建。}

对方式 2 时间戳施加 \(\Delta t^*\)
后，将两组观测合并、排序。若两个观测落在
同一真实时刻，则先检查坐标差，再取均值保留一个节点。对合并后的 6301
个唯一 时刻建立 Cubic Spline，并在

\begin{equation}
t_n=221.0+0.1n,\qquad n=0,1,\ldots,8494
\label{eq:q1-grid}
\end{equation}

上输出 10 Hz 轨迹。输出区间为 \([221.0,1070.4]\)
s，仅覆盖至少一种原始观测 支持的时间区间，不进行外推。

图\ref{fig:q1-objective}表明周期轨迹会形成多个局部低谷，因此互相关只用于候选生成，最终偏差必须由式\eqref{eq:q1-objective}在全可行域内确定。

\begin{figure}[htbp]
  \centering
  \includegraphics[width=.88\textwidth]{q1/objective_scan.pdf}
  \caption{模型一全可行域时间偏差目标函数扫描}
  \label{fig:q1-objective}
\end{figure}

\subsubsection{算法实现}\label{sec:q1-algorithm}

\begin{pseudocodeblock}{算法1\quad 无噪声时间配准与10 Hz轨迹重建}
{\heiti 输入：}两路无噪声定位序列 \(D_1,D_2\)，目标输出频率10 Hz。\par
{\heiti 输出：}时间校正量 \(\Delta t^*\) 与10 Hz二维位置轨迹。\par
\begin{enumerate}[label={\heiti 步骤~\arabic*:},leftmargin=5.5em,itemsep=1pt,topsep=3pt]
  \item 分别构造两路连续轨迹，由速度特征互相关生成时间偏差候选集；
  \item 按式\eqref{eq:q1-overlap}确定可行域，在固定公共网格上计算式\eqref{eq:q1-objective}；
  \item 取全域扫描的最低谷，在其邻域内连续优化，得到式\eqref{eq:q1-optimum}的数值解；
  \item 按式\eqref{eq:q1-time-map}校正方式2时间戳；若两路观测对应同一时刻，则核验坐标后合并；
  \item 对全部唯一时刻建立三次样条，按式\eqref{eq:q1-grid}生成10 Hz轨迹并返回结果。
\end{enumerate}
\end{pseudocodeblock}

\subsubsection{结果分析}\label{sec:q1-results}

采用符号约定
\(t_{2,\mathrm{aligned}}=t_2+\Delta t\)，连续优化得到原始估计

\begin{equation}
\Delta t^*=-198.4316999986\ \mathrm{s}.
\label{eq:q1-raw-offset}
\end{equation}

考虑输入时间戳的显示精度，正式报告为

\begin{equation}
\boxed{\Delta t^*=-198.4317\ \mathrm{s}}.
\label{eq:q1-reported-offset}
\end{equation}

即方式 2 的时间轴需整体向前平移 198.4317 s。对齐公共区间长 700.35 s，
位置 RMSE 为 \(8.51\times10^{-11}\) m。最终 10 Hz 结果共 8495 行。

图\ref{fig:q1-aligned}显示校正后两路观测在公共区间内重合；图\ref{fig:q1-trajectory}给出在原始观测支持域内重建的完整10 Hz轨迹。两图共同说明式\eqref{eq:q1-time-map}只调整时间坐标，没有改变任何原始位置观测。

\begin{figure}[htbp]
  \centering
  \includegraphics[width=.88\textwidth]{q1/aligned_trajectory.pdf}
  \caption{模型一时间校正后的两路定位轨迹（公共区间700.35 s）}
  \label{fig:q1-aligned}
\end{figure}
\begin{figure}[htbp]
  \centering
  \includegraphics[width=.84\textwidth]{q1/trajectory_10hz.pdf}
  \caption{模型一在221.0--1070.4 s内重建的10 Hz完整轨迹}
  \label{fig:q1-trajectory}
\end{figure}

\subsubsection{模型检验与稳健性}\label{sec:q1-validation}

四种插值模型均给出相同的报告偏差。其对齐 RMSE 分别为：Linear
\(4.76\times10^{-6}\) m、Cubic Spline \(8.51\times10^{-11}\) m、PCHIP
\(2.53\times10^{-7}\) m、Akima \(9.91\times10^{-9}\) m，Cubic Spline
最优。

正确偏差下共有 700 组两设备共同采样时刻，最大坐标不一致量为
0，构成不依赖 插值评价网格的独立核验。合成数据测试中，预设偏差
\(-3.7317\) s 的恢复误差为 \(2.20\times10^{-10}\) s。

进一步改变最小重叠时长、粗搜索步长和精细评价步长，共运行 8
个预设情景；原始 时间偏差估计的极差仅为 \(3.56\times10^{-9}\)
s，所有情景按输入精度均报告 \(-198.4317\) s。最优点距搜索边界至少 100.35
s，排除了边界伪最优。

四种插值的误差数量级比较见图\ref{fig:q1-model-comparison}。图中四种方法给出的时间偏差一致，而三次样条的跨设备位置RMSE最低，说明方法选择由无噪声轨迹的局部曲率表达能力决定，并未改变时间偏差结论。

\noindent\textbf{适用边界。}

本模型依赖``固定时间偏差、无空间系统偏差、无随机噪声''三个题面条件。若设备时钟
存在漂移或测量含噪，直接插值会低估不确定性，因此问题二、三需改用偏差辨识和
统计融合。问题一输出的符号约定、公共区间处理和 10 Hz
数据契约将作为后续各问的统一基准。

\begin{figure}[htbp]
  \centering
  \includegraphics[width=.78\textwidth]{paper/model_comparison.pdf}
  \caption{模型一四种插值方法的跨设备位置RMSE比较}
  \label{fig:q1-model-comparison}
\end{figure}
